\scp{Other sound correspondences}{05-04}
PMP *z has a number of reflexes in the \tsc{\ili{Central} Timor} languages
with *z > \it{s} being the most common reflex.
\ili{Welaun} has *z > \it{s,
d}; \ili{Kemak} has *z > \it{s}; \ili{Tokodede} has *z > \it{s,
r}; and \tsc{Mambae} has *z > \it{s} with one possible instance
of *z > \it{r}.
Examples are given in  \trf{05-14}.

\begin{table}\tcp{\tsc{\ili{Central} Timor} *z}{05-14}
%\begin{adjustbox}{width=\textwidth}
	\begin{threeparttable}\st{2.25pt}
	\begin{tabular}{llllllll}\lsptoprule
local gl.	&	path	&	far	&	corn	&	rain	&	ladder	&	bad	&	point\\
PMP	&	*\tbr{z}alan	&	*\tbr{z}auq	&	*\tbr{z}əlay	&	*qu\tbr{z}an	&	*haRə\tbr{z}an	&	*\tbr{z}aqat	&	*tu\tbr{z}uq\\ \midrule
\ili{Welaun}	&	{\hr}\tbr{s}alan	&	{\hr}lai\tbr{s}oo\su{†}	&	{\hr}\tbr{s}ole	&	{\hq}u\tbr{s}an	&	{\hr}o\tbr{s}an	&	{\hr}\tbr{d}aat	&	{\hr}a-tu\tbr{d}u\\
\ili{Kemak}	&	{\hr}\tbr{s}ala	&	\cg{(sə̆goon)}	&	{\hr}\tbr{s}ele	&	{\hq}u\tbr{s}a	&		&	\cg{(mə̆ragan)}	&	\cg{(suku)}\\
\ili{Tokodede}	&	\cg{(daa)}	&	{\hr}k\tbr{r}oo	&	{\hr}\tbr{s}ele	&	{\hq}u\tbr{r}a	&		&	\cg{(klao)}	&	{\hr}t(u)nu\tbr{s}u\\
NW. \ili{Mambae}	&	\cg{(daa)}	&	\cg{(room)}\su{‡}	&	{\hr}\tbr{s}ela	&	{\hq}u\tbr{s}a	&		&	\cg{(klao)}	&	\\
C. \ili{Mambae}	&	\cg{(taan)}	&	{\hr}\tbr{r}oo	&	{\hr}\tbr{s}ela	&	{\hq}u\tbr{s}a	&		&	\cg{(soko)}	&	\\
S. \ili{Mambae}	&	{\hr}\tbr{s}aal	&	{\hr}\tbr{r}oo	&	\cg{(batar)}	&	{\hq}uu\tbr{s}	&		&	\cg{(klao)}	&	\\
		\lspbottomrule
	\end{tabular}
		\begin{tablenotes}\tnsize
			\item[†]	{Initial \it{lai} in \ili{Welaun} \it{laisoo} ‘far’ may be from *daRəq > \it{lai} ‘earth’.}
			\item[‡]	{Final \it{m} in Northwest \ili{Mambae} \it{room} indicates that
								\ili{Central} \ili{Mambae} and South \ili{Mambae} \it{roo} may not be a reflex of *zauq.
								Compare also Northwest \ili{Mambae} from Hatulia \it{roam} ‘far’.}
		\end{tablenotes}
	\end{threeparttable}
%\end{adjustbox}
\end{table}

PMP *b undergoes lenition in all \tsc{\ili{Central} Timor} languages.
\ili{Kemak}, \tsc{Mambae}, \ili{Liquiçá} \ili{Tokodede}
and \ili{Bazartete} \ili{Tokodede} have *b > \it{h} in all positions.
Some varieties of \ili{Kemak} then have sporadic initial *h > {\0}.
\ili{Maubara} \ili{Tokodede} has *b > \it{v} in all positions.
\ili{Welaun} has initial *b > \it{f}
and medial *b > \it{h}.
Examples are given in \trf{05-15}.
Two examples from \ili{Welaun} which are not identical to \ili{Tetun}
and establish the regularity of the pattern are *bəqbəq ‘mouth’
> \it{foha-t} ‘voice’ and *bubuŋ > \it{fuhuk} ‘ridgepole’.

\begin{table}[p]\centering\tcp{\tsc{\ili{Central} Timor} *b}{05-15}
%\begin{adjustbox}{width=\textwidth}
	\st{3.5pt}
	\begin{tabular}{lllllll}\lsptoprule
local gl.	&	pig	&	stone	&	star	&	fruit	&	moon	&	body hair\\
PMP	&	*\tbr{b}a\tbr{b}uy	&	*\tbr{b}atu	&	*\tbr{b}ituqən	&	*\tbr{b}uaq	&	*\tbr{b}ulan	&	*\tbr{b}ulu\\ \midrule
\ili{Welaun}	&	{\hr}\tbr{f}a\tbr{h}i	&	{\hr}\tbr{f}atu	&	{\hr}\tbr{f}itun	&	{\hr}ai \tbr{f}uan	&	{\hr}\tbr{f}ulan	&	{\hr}\tbr{f}ulu-t\\
\ili{Kutubaba} \ili{Kemak}	&		&	{\hr}\tbr{h}atu	&	{\hr}\tbr{h}itu	&	{\hr}\tbr{h}uaŋ	&	{\hr}\tbr{h}ula	&	{\hr}\tbr{h}ulu-ŋ\\
\ili{Kailaku} \ili{Kemak}	&	{\hh}a\tbr{h}i	&	{\hh}atu	&	{\hh}itu	&	{\hr}ai \tbr{h}uan	&	{\hr}\tbr{h}ula	&	{\hr}\tbr{h}ulu-r\\
\ili{Liquiçá} \ili{Tokodede}	&	{\hr}\tbr{h}a\tbr{h}i	&	{\hr}\tbr{h}atu	&	{\hr}\tbr{h}itu	&	{\hr}\tbr{h}uu	&	{\hr}\tbr{h}ula	&	{\hr}\tbr{h}uluː\\
\ili{Maubara} \ili{Tokodede}	&	{\hr}\tbr{v}a\tbr{v}i	&	{\hr}\tbr{v}atu	&	{\hr}\tbr{v}itu	&	{\hr}\tbr{v}uu	&	{\hr}\tbr{v}ula	&	{\hr}\tbr{v}uluː\\
NW. \ili{Mambae}	&	{\hr}\tbr{h}ee\tbr{h}-a	&	{\hr}\tbr{h}ato	&	{\hr}\tbr{h}iut-a	&	{\hr}ai \tbr{h}uan	&	{\hr}\tbr{h}ula	&	{\hr}\tbr{h}ulu-n\\
C. \ili{Mambae}	&	{\hr}\tbr{h}ae\tbr{h}-a	&	{\hr}\tbr{h}aut-a	&	{\hr}\tbr{h}iut-a	&	{\hr}ai \tbr{h}uan	&	{\hr}\tbr{h}ula	&	{\hr}\tbr{h}ulu-n\\
S. \ili{Mambae}	&	{\hr}\tbr{h}ae\tbr{h}	&	{\hr}\tbr{h}aut	&	{\hr}\tbr{h}iut	&	{\hr}ai \tbr{h}ua	&	{\hr}\tbr{h}ulai	&	{\hr}\tbr{h}ulu\\
		\lspbottomrule
	\end{tabular}
%\end{adjustbox}
\end{table}

Word initially PMP *k > {\0} in all \tsc{\ili{Central} Timor} languages
except \ili{Tokodede} where it is retained as \it{k}.
Word medially *k > \it{ʔ} in \ili{Welaun}; *k > \it{ʔ,} {\0} in \ili{Kemak};
*k > \it{ʔ, k} in \ili{Tokodede}; and *k > {\0} in most \tsc{Mambae} words,
though South \ili{Mambae} retains *k = \it{k} in at l\ili{east} one word.
Examples of PMP *k are given in \trf{05-16}.

\begin{table}[p]\centering\tcp{\tsc{\ili{Central} Timor} *k}{05-16}
\begin{adjustbox}{width=\textwidth}
	\begin{threeparttable}\st{2pt}
	\begin{tabular}{llllllll}\lsptoprule
local gl.	&	wood, tree	&	\tsc{2sg}	&	head louse	&	tail	&	steal	&	ascend	&	elbow\\
PMP	&	*\tbr{k}ahiw	&	*\tbr{k}ahu	&	*\tbr{k}utu	&	*i\tbr{k}uR	&	*na\tbr{k}aw	&	*sa\tbr{k}ay	&	*si\tbr{k}u\\ \midrule
\ili{Welaun}	&	{\hr}ai	&	{\hr}oo	&	{\hr}utu	&	{\hr}hi\tbr{ʔ}u-n\su{†}	&	{\hr}a-na\tbr{ʔ}o	&	{\hr}sa\tbr{ʔ}e	&	{\hr}si\tbr{ʔ}u-t\\
\ili{Kemak}	&	{\hr}ai	&	{\hr}oo	&	{\hr}utu	&	{\hr}hi\tbr{ʔ}o-r\su{†}	&	{\hr}pə̆nao	&	{\hr}sae	&	{\hr}si\tbr{ʔ}u-r\\
\ili{Tokodede}	&	{\hr}\tbr{k}ai	&	{\hr}\tbr{k}oo	&	{\hr}\tbr{k}utu	&	{\hr}i\tbr{k}o	&	{\hr}mana\tbr{ʔ}o	&	{\hr}sa\tbr{ʔ}e	&	\\
NW. \ili{Mambae}	&	{\hr}ai-a	&		&	{\hr}uut-a	&	{\hr}io-n	&	{\hr}pnao	&	{\hr}leol sae-n\su{‡}	&	{\hr}siu-n\\
C. \ili{Mambae}	&	{\hr}ai-a	&	{\hr}oo	&	{\hr}utu-n	&	{\hr}io-n	&	{\hr}fnao	&	{\hr}leel-a sae\su{‡}	&	{\hr}lima siu-n\\
S. \ili{Mambae}	&	{\hr}ai	&	{\hr}oo	&	{\hr}uut	&	{\hr}io	&	{\hr}fnao-tee	&	{\hr}leol sae\su{‡}	&	{\hr}lima si\tbr{k}u\\
		\lspbottomrule
	\end{tabular}
		\begin{tablenotes}\tnsize
			\item[†]	{\ili{Welaun} \it{hiʔu-n}
								and \ili{Tokodede} \it{hiʔo-r} ‘tail’ have insertion of /h/ /{\#\gap}VʔV.}
			\item[‡]	{\tsc{Mambae} forms given here mean ‘\ili{east}’, literally ‘sun rise’.}
		\end{tablenotes}
	\end{threeparttable}
\end{adjustbox}
\end{table}

The usual reflexes of word-final vowel-glide sequences are *-ay
> \it{e} and *-aw > \it{o},
though there are also sporadic reflexes of *-aw/-*ay > \it{a} in \tsc{Mambae}.
Additionally, all \tsc{Nuclear \ili{Central} Timor} languages have
*-aw > \it{a} in reflexes of *lakaw ‘go’,
with \ili{Welaun} showing *-aw > \it{a {\tl} o} in this word.
Examples are given in \trf{05-17}.

\begin{table}[p]\centering\tcp{\tsc{\ili{Central} Timor} *-ay and *-aw}{05-17}
\begin{adjustbox}{width=\textwidth}
	\st{2.5pt}
	\begin{tabular}{lllll|llll}\lsptoprule
local gl.	&	liver	&	die	&	snake	&	corn	&	steal	&	sun	&	rat	&	go\\
PMP	&	*qat\tbr{ay}	&	*mat\tbr{ay}	&	*nip\tbr{ay}	&	*zəl\tbr{ay}	&	*nak\tbr{aw}	&	*qaləj\tbr{aw}	&	*balab\tbr{aw}	&	*lak\tbr{aw}\\ \midrule
\ili{Welaun}	&	{\hr}at\tbr{e}-t	&	{\hr}mat\tbr{e}	&	\cg{(saba)}	&	{\hr}sol\tbr{e}	&	{\hr}a-naʔ\tbr{o}	&	{\hr}lor\tbr{o}	&	{\hr}lah\tbr{o}	&	{\hr}laʔ\tbr{a}, laʔ\tbr{o}\\
\ili{Kemak}	&	{\hr}at\tbr{e}-r	&	{\hr}mat\tbr{e}	&	{\hr}nip\tbr{e}	&	{\hr}sel\tbr{e}	&	{\hr}pə̆na\tbr{o}	&	{\hr}lel\tbr{o}	&	{\hr}lah\tbr{o}	&	{\hr}la\tbr{a}\\
\ili{Tokodede}	&	{\hr}at\tbr{e}	&	{\hr}mat\tbr{e}	&	{\hr}nip\tbr{e}	&	{\hr}sel\tbr{e}	&	{\hr}manaʔ\tbr{o}	&	{\hr}lel\tbr{o}	&		&	{\hr}la\tbr{a}\\
NW. Mam.	&	{\hr}at\tbr{e}-n	&	{\hr}me\tbr{e}t	&	\cg{(sokai-a)}	&	{\hr}sel\tbr{a}	&	{\hr}pna\tbr{o}	&	{\hr}lel\tbr{o}n	&	{\hr}lah\tbr{o}	&	\\
C. \ili{Mambae}	&	{\hr}at\tbr{e}-n	&	{\hr}ma\tbr{e}t	&	\cg{(solai-a)}	&	{\hr}sel\tbr{a}	&	{\hr}fna\tbr{o}	&	{\hr}leel-a	&	{\hr}lah\tbr{a}	&	\\
S. \ili{Mambae}	&	{\hr}at\tbr{e}	&	{\hr}ma\tbr{e}t	&	\cg{(sohai)}	&	\cg{(batar)}	&	{\hr}fna\tbr{o}-tee	&	{\hr}le\tbr{o}l mata	&	{\hr}la\tbr{o}h	&	{\hr}la\tbr{a}\\
		\lspbottomrule
	\end{tabular}
\end{adjustbox}
\end{table}

PMP *w is usually retained unchanged word initially in \ili{Welaun}
with one example of *w > \it{b} medially in *sawa > \it{saba} ‘snake’.
PMP *w is also usually lost in \ili{Tokodede} and \tsc{Mambae}.
\ili{Kemak} usually has *w > \it{b},
though *wakaR ‘roots’ > \it{haʔa-n} has *w > {\0} with subsequent
epenthesis of \it{h} /{\#}{\gap}VʔV (compare *ikuR > \it{hiʔo-r} ‘tail’).
Examples are given in \trf{05-18}.

\begin{table}\tcp{\tsc{\ili{Central} Timor} *w}{05-18}
%\begin{adjustbox}{width=\textwidth}
	\begin{threeparttable} \st{3.5pt}
	\begin{tabular}{llllll}\lsptoprule
local gl.	&	water	&	roots	&	eight	&	day	&	nine\\
PMP	&	*\tbr{w}ahiyəR	&	*\tbr{w}akaR	&	*\tbr{w}alu	&	*\tbr{w}aRi\su{Δ}	&	*si\tbr{w}a\\ \midrule
\ili{Welaun}	&	{\hr}\tbr{w}ee	&	{\hr}\tbr{w}aʔa-n	&	{\hr}\tbr{w}alu	&	{\hr}\tbr{w}ainrua	&	{\hr}si\tbr{w}i\\
\ili{Kemak}	&	{\hr}\tbr{b}ia, \tbr{b}ea\su{†}	&	{\hr}haʔa-n\su{‡}	&	{\hr}\tbr{b}alu	&	{\hr}\tbr{b}airua	&	{\hr}si\tbr{b}e\\
\ili{Tokodede}	&	{\hr}ee	&	\cg{(ramut)}	&	\cg{(hoho telu)}	&	{\hr}airuu-n	&	\cg{(hoho paat)}\\
NW. \ili{Mambae}	&	{\hr}eer-a	&	\cg{(ramu-n)}	&	\cg{(hoho teul)}	&	{\hr}airua-n	&	\cg{(hoho paat)}\\
C. \ili{Mambae}	&	{\hr}eer-a	&	\cg{(ramu-n)}	&	{\hr}\tbr{u}alu\su{ǁ}	&	{\hr}airua-n	&	{\hr}sia\su{ǁ}\\
S. \ili{Mambae}	&	{\hr}eer	&	{\hr}haa	&	\cg{(liim nai telu)}	&	{\hr}airua	&	\cg{(liim nai fata)}\\
		\lspbottomrule
	\end{tabular}
		\begin{tablenotes}\tnsize
			\item[†]	{Most varieties of \ili{Kemak} have \it{bia} ‘water’.
								\ili{Marobo} and \ili{Leosibe} \ili{Kemak} have \it{bia {\tl} bea}.}
			\item[‡]	{\ili{Kemak} \it{haʔa-n} ‘roots’ has insertion of /h/ /{\#\gap}VʔV.
								South \ili{Mambae} \it{haa} ‘roots’ probably also shows this process
								with subsequent medial **ʔ > {\0}.}
			\item[Δ] 	{See note to \trf{05-05} (\prf{tab:05-05}) for the glosses
								of reflexes of *waRi given here.}
			\item[{ǁ}]	{\ili{Central} \ili{Mambae} \it{ualu} ‘eight’
								and \it{sia} ‘nine’ may be borrowings from eastern \ili{Tetun}.}
		\end{tablenotes}
	\end{threeparttable}
%\end{adjustbox}
\end{table}

\largerpage
Some varieties of \ili{Kemak},
including \ili{Saneri} \ili{Kemak} and \ili{Marobo} \ili{Kemak} have *t > \it{ʧ} before high vowels.
\ili{Saneri} \ili{Kemak} also appears to have *t > \it{ʧ} word finally after
high vowels, though there is only one example in our data.
Examples are given in \trf{05-19}.

\begin{table}
	\centering\tcp{\ili{Kemak} *t > \it{ʧ} /V\tsc{[+high]}}{05-19}
	\st{2.25pt}\begin{tabular}{llllllll}\lsptoprule
local gl.			&	head louse				&	roast						&	skin									&	white						&	stone						&	old								&	grass\\
PMP						&	*ku\tbr{t}u				&	*\tbr{t}unu			&	*kuli\tbr{t}					&	*ma-pu\tbr{t}iq	&	*ba\tbr{t}u			&										&	\\ \midrule
\ili{Kailaku} \ili{Kemak}	&	\hp{*k}u\tbr{t}u	&									&	\hp{*k}uli\tbr{t}-iir	&	{\hr}bu\tbr{t}i	&	\hp{*b}a\tbr{t}u&	{\hr}\tbr{t}uman	&	{\hr}suu\tbr{t}\\
\ili{Marobo} \ili{Kemak}	&	\hp{*k}u\tbr{ʧ}u	&	{\hr}\tbr{ʧ}unu	&	\hp{*k}uli\tbr{ʧ}-iir	&	{\hr}bu\tbr{ʧ}i	&	\hp{*b}a\tbr{ʧ}u&	{\hr}\tbr{ʧ}uman	&	{\hr}suu\tbr{t}\\
\ili{Saneri} \ili{Kemak}	&	\hp{*k}u\tbr{ʧ}u-r&	{\hr}\tbr{ʧ}unu	&	\hp{*k}ul\tbr{ʧ}i-r		&	{\hr}bu\tbr{ʧ}i	&	\hp{*b}a\tbr{ʧ}u&	{\hr}\tbr{ʧ}uman	&	{\hr}suu\tbr{ʧ}\\
		\lspbottomrule
	\end{tabular}
\end{table}

Finally, \ili{Tokodede} and some varieties of \ili{Kemak} have loss of final
*n with compensatory lengthening of the previous vowel.
Some other varieties of \ili{Kemak} have *n > \it{ŋ} syllable finally,
and this may represent an intermediate stage before loss of the nasal.
Thus, the full pathway may have been *Vn > **Vŋ > **Vː.
This process is attested in \ili{Tokodede}, \ili{Marobo} \ili{Kemak},
\ili{Lemia} \ili{Kemak}, \ili{Diirbatin} \ili{Kemak}, and \ili{Atsabe} \ili{Kemak}.
Examples are given in \trf{05-20}.

\begin{table}
	\centering\tcp{\ili{Kemak} and \ili{Tokodede} *n > {\0} and compensatory lengthening}{05-20}
	\begin{adjustbox}{width=\textwidth}
	\st{2pt}\begin{tabular}{lllllll}\lsptoprule
local gl.	&	new	&	hot	&	body hair	&	bark	&	straight	&	big\\
PMP	&	*baqəRu	&	*ma-panas	&	*bulu	&	*kayu + *kulit	&		&	\\ \midrule
\ili{Leosibe} \ili{Kemak}	&	{\hr}heu\tbr{n}	&	{\hr}bansa\tbr{n}	&	{\hr}hulu-r	&	{\hr}ai ulti-\tbr{n}	&	{\hr}loso\tbr{n}	&	{\hr}bote\tbr{n}\\
\ili{Leolima} \ili{Kemak}	&	{\hr}heu\tbr{ŋ}	&	{\hr}baŋsa\tbr{ŋ}	&	{\hr}hulu-\tbr{ŋ}	&	{\hr}au ulti\tbr{ŋ}	&	{\hr}loso\tbr{ŋ}	&	{\hr}bote\tbr{ŋ}\\
\ili{Lemia} \ili{Kemak}	&	{\hr}heu\tbr{ː}	&	{\hr}bansa\tbr{ː}	&	{\hr}hulu\tbr{ː}	&	{\hr}ai ulit	&	{\hr}loso\tbr{ː}	&	{\hr}bote\tbr{ː}\\
\ili{Diirbatin} \ili{Kemak}	&	{\hr}heu\tbr{ː}	&	{\hr}bansa\tbr{ː}	&	{\hr}lima-r hulu\tbr{ː}	&	{\hr}ai ulti\tbr{ː}	&	{\hr}loso\tbr{ː}	&	{\hr}bote\tbr{ː}\\
\ili{Tokodede}	&	{\hr}heu\tbr{ː}	&	{\hr}bana\tbr{ː}	&	{\hr}hulu\tbr{ː}	&	{\hr}kuluta\tbr{ː}	&		&	\cg{(aʔe)}\\
		\lspbottomrule
	\end{tabular}
	\end{adjustbox}
\end{table}

